%% 執筆用サンプル（main.tex からは読み込まれません）
%% 各章で使える環境・マクロの書き方例です。

\chapter{章タイトル}
\label{chap:example}

\begin{goalbox}
  \begin{itemize}
    \item 学ぶこと 1
    \item 学ぶこと 2
  \end{itemize}
\end{goalbox}

\section{節タイトル}

インラインでコマンドを書くときは \cmd{ls -l} のようにします。
キー操作は \key{Ctrl}+\key{C} のように書きます。

% ターミナル操作例（タイトルは省略可）
\begin{terminal}[ファイル一覧を表示する]
$ ls -l ~/ros2_ws/src
total 4
drwxrwxr-x 3 user user 4096 Sep 29 18:33 my_package
\end{terminal}

% ソースコード
\begin{codelisting}[style=python]{最小のノード}{lst:minimal}
import rclpy
from rclpy.node import Node

def main():
    rclpy.init()
    node = Node('minimal')
    rclpy.spin(node)
    rclpy.shutdown()
\end{codelisting}

% 外部ファイルのコードを読み込む場合
% \lstinputlisting[style=python]{samples/rclpy/my_package/my_package/talker.py}

\begin{point}
  覚えておきたいポイント。
\end{point}

\begin{caution}[rm コマンドに注意]
  \cmd{rm -rf} は取り消せません。
\end{caution}

\begin{column}{ROS 1 との違い}
  コラム本文。
\end{column}

% 図
% \begin{figure}[htbp]
%   \centering
%   \includegraphics[width=0.8\linewidth]{ros2_concepts/rqt_graph.png}
%   \caption{rqt\_graph の表示例}
%   \label{fig:rqt_graph}
% \end{figure}

\todo{あとで書く}

\section*{この章のまとめ}
